%%%%%%%%%%%%%%%%%%%%%%%%%%%%%%%%%%%%%%%%%%%%%%%%%%%%%%%%%%%%
% Pedro Brandão's trial to get a template for thesis for students
% Used the upthesis from Fernando Silva (see upthesis).
% See also the packages file.
% 2014/07/07 First draft
% 2014/07/21 
%  pbrandao: added the list of listings (it should produce portuguese name if
%           babel is set to portuguese, see packages.tex). Changed usepackage of babel
%           to be before input packages.tex to allow test

%
%%%%%%%%%%%%%%%%%%%%%%%%%%%%%%%%%%%%%%%%%%%%%%%%%%%%%%%%%%%%


% makes all pages the height of the text on that page. No extra vertical space is added.
\raggedbottom 
% setting it to report will remove the blank pages before each chapter
\documentclass[11pt,a4paper,twoside]{book}


%%%%%%%%%%%%%%%%%%%%%%%%%%%%%%%%%%%%%%%%%%%%%%%%%%%%%%%%%%%%
%%%   Packages that need to be configured for the thesis
%%%%%%%%%%%%%%%%%%%%%%%%%%%%%%%%%%%%%%%%%%%%%%%%%%%%%%%%%%%%

% Language settings
% use UKenglish for UK or leave blank for US English
% it will also change the names for some of the chapter (list of tables, figures, content,
%\usepackage[UKenglish]{babel}
\usepackage[portuguese]{babel}

%%%%%%%%%%%%%%%%%%%%%%%%%%%%%%%%%%%%%%%%%%%%%%%%%%%%%%%%%%%%
%%%   Packages uses language definitions
% see file below for more packages and settings
%%%%%%%%%%%%%%%%%%%%%%%%%%%%%%%%%%%%%%%%%%%%%%%%%%%%%%%%%%%%
%%%%%%%%%%%%%%%%%%%%%%%%%%%%%%%%%%%%%%%%%%%%%%%%%%%%%%%%%%%%
% Pedro Brandão adaptation from UPThesis of Fernando Silva
% (see upthesis.sty)
% Added more packages and configuration
% Should be used with pdflatex for better output
%
% 2014/07/08 First draft
% 2014/07/21 added the List of code blocks and their ``name'' for portuguese
%        this includes the caption name
%
% TODO there should be many...
% example add the required cover page: current solution is
% to print it using word to pdf and then join the pdfs
%%%%%%%%%%%%%%%%%%%%%%%%%%%%%%%%%%%%%%%%%%%%%%%%%%%%%%%%%%%%

% Particular settings for the layout of the thesis, including the old FC logo
% The header format is also on that file
\usepackage{upthesis}

% Encoding
\usepackage[utf8]{inputenc}
% see http://tex.stackexchange.com/questions/664/why-should-i-use-usepackaget1fontenc
\usepackage[T1]{fontenc}

%image placement
\usepackage{float}
\usepackage{graphicx}

% Use colour names
\usepackage[usenames,dvipsnames]{color}
% for named colors (see 
% http://en.wikibooks.org/wiki/LaTeX/Colors#The_68_standard_colors_known_to_dvips)
\definecolor{light-gray}{gray}{0.94}
\definecolor{gray-cell}{gray}{.80}
\usepackage{colortbl}
\usepackage{tabularx}
\usepackage{multicol}
\usepackage{multirow}


% rotating introduces some glitches with acronym. Acronyms used in the sidewaystable were marked as being multiply
% defined. They were expanded the 1st time on the table. Using \acs solved the issue
\usepackage{rotating}

% used for code examples and other listings
\usepackage{listings}
\lstset{
  aboveskip={\baselineskip},
  % belowskip={1.5\baselineskip},
 % columns=fixed,
 % extendedchars=true,
  % prebreak = \raisebox{0ex}[0ex][0ex]{\ensuremath{\hookleftarrow}},
  %identifierstyle=\bf\ttfamily,
  keywordstyle=\color{Maroon}\bf\ttfamily,
  commentstyle=\normalfont\color{BlueViolet}\itshape,
  stringstyle=\color{OliveGreen},
%  lineskip=-2.0pt,
  %from http://en.wikibooks.org/wiki/LaTeX/Packages/Listings
%  language=Octave,                % choose the language of the code
  basicstyle=\footnotesize\ttfamily,       % the size of the fonts that are used for the code
%  numbers=left,                   % where to put the line-numbers
%  numberstyle=\footnotesize,      % the size of the fonts that are used for the line-numbers
%  stepnumber=2,                   % the step between two line-numbers. If it's 1 each line 
%                                  % will be numbered
  numbersep=5pt,                  % how far the line-numbers are from the code
  backgroundcolor=\color{light-gray},  % choose the background color. You must 
                                       %add \usepackage{color}
%  showspaces=false,               % show spaces adding particular underscores
  showstringspaces=false,         % underline spaces within strings
%  showtabs=false,                 % show tabs within strings adding particular underscores
%  frame=shadowbox,
  frame=single,	                % adds a frame around the code
  frameround=tttt,	                % t or f(alse) for each corner
  framesep = 1pt,
  %framexleftmargin = -10pt,
  fillcolor=\color{white},
  rulecolor=\color{MidnightBlue},
  %rulesepcolor=\color{blue},
  tabsize=3,	                % sets default tabsize to 3 spaces
  captionpos=b,                   % sets the caption-position to bottom
  xleftmargin = 12pt,
  xrightmargin = 10pt,
  breaklines=true,                % sets automatic line breaking
  breakatwhitespace=true,        % sets if automatic breaks should only happen at whitespace
%  title=\lstname,                 % show the filename of files included with \lstinputlisting;
%                                  % also try caption instead of title
%  escapeinside={\%*}{*)},         % if you want to add a comment within your code
%  morekeywords={*,...}            % if you want to add more keywords to the set
   escapechar=\#,                   % enable to right latex and be interpreted as such
}

\makeatletter
  \@ifpackagewith{babel}{portuguese}{%
    \addto\captionsportuguese{\renewcommand*\lstlistlistingname{Lista de Blocos de Código}}%
    \addto\captionsportuguese{\renewcommand*\lstlistingname{Bloco de Código}}%
  }{}
\makeatother

% This adds the possibility of having lists in paragraphs
\usepackage{paralist}

% Enhance/improve fonts, kerning and spacing
\usepackage{microtype}
% Need to be a font compatible with microtype (or remove microtype package)
\usepackage{lmodern} %% See for fonts http://www.tug.dk/FontCatalogue/ and http://www.tug.org/fonts/
%fonts examples: newcent, utopia, charter, concmath (may need cm-super)
% can also add the cm-super package (don't usepackage on the tex file,
% just install using the package manager)

% for monotype use courier
% http://www.macfreek.nl/mindmaster/LaTeX_Bold_Typewriter_Font
\usepackage{courier}
% fix the restriction for font-sizes with latex
% see http://www.tex.ac.uk/cgi-bin/texfaq2html?label=fontunavail
\usepackage{fix-cm}

% Math fonts
\usepackage{amsfonts}
\usepackage{amsmath}

% Bibliography
% for natbib see http://en.wikibooks.org/wiki/LaTeX/More_Bibliographies#Options
%\usepackage[authoryear,sort]{natbib}
\usepackage[numbers,sort&compress]{natbib}

% Acronyms package
\usepackage{acronym}

% Enable multicoluns, used in the acronyms appendix
\usepackage{multicol}

% just for testing purposes (get ipsum text)
\usepackage{lipsum}

% URLs, linking, and pdf parameters
% break long urls, should be before hyperref
\usepackage[hyphens]{url}

%%%%%%%%%%%%%%%%%%%%%%%%%%%%%%%%%%%%%%%%%%%%%%%%%%
%%% Headers/Footers/Chapters
%%%%%%%%%%%%%%%%%%%%%%%%%%%%%%%%%%%%%%%%%%%%%%%%%%
% See the upthesis for that


\usepackage[
%backref={section},
%pagebackref, % for getting references to the page where the citation is (in the biblio)
pdfpagelabels=false
]{hyperref}
\hypersetup{pdftitle={Titulo da tese}, %nao suporta acentos
   pdfkeywords ={palavras chave},
   pdfsubject = {assunto},
	bookmarksnumbered=true,
   pdfauthor ={Autor}, % see other options on manual (can be page) needs empty line on bibitem
   plainpages=false, 
   pdfborder={0 0 0},
   colorlinks,%colorlinks=false,
   breaklinks=true,
	%linktocpage= false, make page number, not text, be link on TOC, LOF and LOT 
	%hyperindex=true 	% Makes the page numbers of index entries into hyperlinks. Relays on unique page anchors (pageanchor) 
% see for colors http://mirror.ctan.org/macros/latex/contrib/xcolor/xcolor.pdf
   linkcolor=	Sepia, %MidnightBlue,% BlueViolet,%Sepia, % Color for normal internal links.
   %anchorcolor=black,% Color for anchor text.
   citecolor=RedViolet,% Color for bibliographical citations in text.
   %filecolor=cyan% Color for URLs which open local files.
   %menucolor=red% Color for Acrobat menu items.
   %runcolor=filecolor% Color for run links (launch annotations).
   urlcolor=NavyBlue% Color for linked URLs.
}
%use the same style for \url as the text
% from http://en.wikibooks.org/wiki/LaTeX/Hyperlinks#Customization
\urlstyle{same}


\begin{document}

\title{Título da Tese}
\submitionplace{Tese submetida à Faculdade de Ciências da \\
  Universidade do Porto para obtenção do grau de Mestre \\ 
  em Ciência de Computadores}
\author{Nome do autor}
\department{Departamento de Ciência de Computadores \\ Faculdade de
  Ciências da Universidade do Porto}
\submitdate{Setembro 2015}


\beforepreface%



\prefacesection{Abstract}

A long, long time ago... 

\prefacesection{Resumo}
Há muito, muito tempo


\prefacesection{Agradecimentos}

Obrigado a todos, obrigado \ldots


\dedicationpage{Dedico a \ldots}

% end of thesis preamble
\afterpreface%

%% main tex here
%% By putting the chapter names here, one can just comment the content in the chapters 
%% and produce a pdf with the correct chapter number.
%% If you want further configurability you can use subfiles package
%% https://www.ctan.org/pkg/subfiles

%\chapter{Introdução}\label{chap:intro}
\chapter{Introdução}\label{chap:intro}

Este template pretende dar uma base para a escrita de dissertações no \ac{DCC}/\ac{FCUP}.

\section{Organização}
Vejam o capítulo~\ref{chap:back} para algum contexto sobre o conteúdo de uma dissertação. O capítulo~\ref{chap:devel} tem alguns pontos sobre o uso do {\LaTeX} e ao longo do texto são utilizados alguns exemplos para mostrar algumas das funcionalidades.






%\chapter{Background}\label{chap:back}
\chapter{Background}\label{chap:back}

This chapter has the ``Instructions for preparing and writing M.Sc. Dissertations'' (Research-oriented work), 
Version 1.0, January, 7$^{th}$, 2015 from prof. Inês Dutra. It was originally written in English so it was kept as such. It was some additions from Pedro Brandão.

\section{Before Starting}
%%%%%%%%%%%%%%%%%%%%%%%%%%%%%%%%%%%%%%%%
Before starting your dissertation, you need to \textbf{define} what is the subject you are going to work
with and perform a \textbf{thorough systematic} bibliography review of the theme you chose.
What is a systematic review? It is the one where you search the web or books for the subject, and
define rules for filtering papers in two sets “included” and “excluded” and explain why some papers
go to one set or the other.
In order to start the search, you need to prepare keywords related to your subject and prepare
queries to be used in Google Scholar, Scopus, MesH etc. These search engines will return a number
of papers on the subject you are looking for.\textbf{You need to read at least the abstract and 
conclusions of every paper retrieved after your search}. Now, you filter out only the ones you
think are very closely related to your research work, and give a reason for choosing those papers.

\section{Bibliography}
%%%%%%%%%%%%%%%%%%%%%%%%%%%%%%%%%%%%%%%%
 Start organizing your bibliography file. Choose one of the standards available to start organizing
your references. Usually your department/faculty has clear rules about the standard to be used. If
you are formatting your text using \LaTeX, most references found during your bibliographic search
can be exported in BibTeX format.

\subsection{Bibliography Section}
%---------------------------------------
Your dissertation needs to have a Bibliography Section with a list of the cited works you have in
the text. In the process of writing your dissertation, make sure to properly refer the authors you are
basing your text on. For example, “Yaacoub et al.~\cite{yaacoub2012} discovered that\ldots”. In this sentence, Yaacoub is the
first author of one of the publications you list in the Bibliography Section of your dissertation and~\cite{yaacoub2012} is the link that connects this citation to the publication in the bibliographic list. If your
bibliographic entry has only one author, you cite only the author's surname. If the bibliographic
entry has two authors you cite the two authors' surnames (e.g., Clausen and Jacquet~\cite{Clausen2003}). If the
bibliographic entry has more than two authors, you can use the expression “et al.”, like in the
example shown before.

You should avoid as much as possible web site references. Only in some cases, illustration, showing trends, are they acceptable.

References should be used to back up claims made. Specially in the introduction section, sentences that stipulate something should be backed by references that assert that claim (e.g.: Android, in December 2016, was the most widely used mobile operating system~\cite{netMarketShareMobileOS}\foonote{An example where a web link to a recognized market analysis company would be valid. Note that the time which the report was seen is very relevant.)

\section{Inserts}
%%%%%%%%%%%%%%%%%%%%%%%%%%%%%%%%%%%%%%%%
Every picture, graph, diagram, algorithm etc needs to have a caption and a number and needs to
be cited and explained in the text. The mere existence of a picture, etc. does not exempt it of a description. 

\subsection{Copyrights and image usage}
%---------------------------------------
If you want to use any picture, graph, diagram etc available in one of the publications in your
dissertation, you need to make sure that you can use it (check the copyright rules). If the copyright
rules allow you to reproduce the picture (or others) in your text, you need to insert a reference to the
source (where the picture was taken from) in the caption. If you are allowed to use a picture, but
want to slightly modify it, you need to say in the caption: Adapted from [1] (where [1] is the
number of your reference in the Bibliographic Section).


\section{Chapters/Organization}
%%%%%%%%%%%%%%%%%%%%%%%%%%%%%%%%%%%%%%%%
\begin{description}
   \item[Chapter 1: Introduction:]
This should be a summary of what comes in the next chapters. Here you explain in two to five
pages: (1) the context of your work highlighting and defining the problem you need to solve, (2)
what you want to do (objectives) and (2) why you want to do it (motivation), (3) how you want to
achieve your objectives (methodology), always supporting your text on the available literature, (4)
contributions (Results that confirm that you achieved your objectives) and (5) organization of the
chapters that come next.
\item[Chapter 2: Basic Concepts:]
In this chapter you need to present the foundations of your work: theoretical aspects, background
material etc, all that is needed to understand the terminology and expressions used in the remaining
chapters.
\item[Chapter 3: Related Work:]
Here you need to discuss about other works in the literature that do something similar to what
you want to do. You need to cite and discuss the relevant papers you chose to include in your study
during your survey. Explain what others do, why it is not sufficient, and why you need to do what
you want to do. It is helpful to define some criteria to compare your work against others, and
build a table with main characteristics of other works contrasting to what you want to do. In other
words, in which aspects is your work different from others?
\item[Chapter 4: Your Work:] this chapter describes the contributions of the work done. If it is based on prior work (continuation of the project or using prior developed work), the should only describe what is the new work done. If references are needed, it should be clear what is the prior work and what is the new contribution.
\item[Chapter 5: Materials and Methods:] should have:
   \begin{itemize}
      \item Definition of Experiments (if any)
      \item Definition of Evaluation Metrics
   \end{itemize}
\item[Chapter 6: Results and Analysis]
\item[Chapter 7: Conclusions and Future work:] 
         this should restate the problem and iterate through the solution(s) analyzing the advantages and contributions. The limitations and unsolved problems should also be described.
         It should also describe the potentiality of new research/development that the work enables, the future work.
   \begin{itemize}
      \item Research Summary
      \item Main Findings
      \item Limitations
      \item Future Work
      \item Conclusion
   \end{itemize}
\end{description}
During writing, some of these chapters may collapse into just one.\textbf{Your work (chapters 4-7) should account for at least 50\% of your whole dissertation.}

\subsection{Contents of each chapter}
%---------------------------------------
You should start each chapter with a summary of it is its objective and contents, preferably relating to previous ones. At the end of the chapter provide a conclusion/summary of it, preferably connecting it to the next one.

\section{Cover}
%%%%%%%%%%%%%%%%%%%%%%%%%%%%%%%%%%%%%%%%
 Your dissertation needs to have a cover, and a table of contents. It is also useful to have, in the
preamble, a list of Figures, Tables, Algorithms etc. It is recommended to have a look at other recent
dissertations of your colleagues of the same department to make sure you use the right formatting
rules. These dissertations are publicly available online in the digital repository of the University of
Porto. For example, the FCUP repository can be found at:
\url{https://repositorio-aberto.up.pt/handle/10216/9535}


%\chapter{Estado da Arte}\label{chap:stat}
\chapter{Estado da Arte}\label{chap:stat}

\section{Section example}
\lipsum[1-6]
\section{Second Section example}
\subsection{SubSection example}
\lipsum[7-10]


\begin{table}[htbp]
   \caption[BSN vs WSN]{BSN versus WSN
   (with input from  Latré  and Guang )}
\label{tab:bsnvswsn}
\centering
{
\footnotesize
\begin{tabularx}{0.98\textwidth}{|>{\columncolor{gray-cell}}c|X|X|}
   \hline
   \rowcolor{gray-cell} 
   &    \centering  \textbf{BSN} &  \centering \textbf{WSN}  \tabularnewline 
   \hline
   %%%% line
   \begin{sideways} \hspace{-11em} \textbf{Distribution} \end{sideways}
   &  
   \begin{asparaenum}[\bfseries i)]
      \item Existence of a \ac{BS};
      \item \ac{BS} collects, maintains and processes the data;
      \item Nodes will do minimal processing, sending all data to the \ac{BS};
      \item Centralized system where \ac{BS} controls all nodes;
      \item Node replacement is difficult in in-body sensor nodes;
      \item Smaller number of nodes;
      \item Nodes need to take biocompatibility, wearability into account.
   \end{asparaenum}
   & 
   \begin{asparaenum}[\bfseries i)]

      \item A \ac{BS} may or not exist or there may be several \acp{BS} (e.g. mobile nodes 
         collect info, clustering);
      \item As in \ac{BSN}, but also on-demand querying;
      \item Nodes will do processing, aggregation to alleviate communication or 
         correlate results;
      \item Distributed system, nodes decide cooperatively;
         %(form clusters, aggregate data, \etc);
      \item Node replacement is difficult due to location, scale, etc.;
      \item (usually) Wide areas covered by large number of nodes.
      \item Nodes may need to be environment friendly, indiscernible from surroundings.
     %    \vspace{-0.7em}
   \end{asparaenum}
   \tabularnewline \hline
   %%%% line
   \textbf{Comm.} 
   & 
   \begin{asparaenum}[\bfseries i)]
      \item One hop to \ac{BS};
      \item Close range but attenuated by body;
      \item Data rates heterogeneous.
    %\vspace{-0.7em}
   \end{asparaenum}
   & 
   \begin{asparaenum}[\bfseries i)]
      \item Multi hop through network of sensor nodes;
      \item Long(er) range;
      \item Data rates homogeneous.
    %\vspace{-0.7em}
   \end{asparaenum}
   \tabularnewline \hline
   %%%% line
    \textbf{Data} 
   & 
   This is some text on this cell. The multirow package does not know the height of the cell and can not center the cell to the right. This is because of the X from tabularx.
   & 
   \multirow{ 2}{*}{This is on two rows of the table}

 \tabularnewline \cline{1-2}\noalign{\vskip.3pt}% using the noaling vskip to show the line
   %%%% line
    \textbf{Energy} 
   & 
   Some more text just to show something
   %      \vspace{-0.7em}
   & 
    \tabularnewline \hline
\end{tabularx}
}
\end{table}

%% vim: set tw=80 spl=en_gb spell: syntax spell toplevel  :



\lipsum[4-9]







%\chapter{Desenho da Arquitetura}\label{chap:syst}
\chapter{Desenho e Desenvolvimento}\label{chap:devel}

O referente ao ``Your work'' do capítulo~\ref{chap:back}.

No contexto deste template, este capítulo serve de exemplos de uso do \LaTeX e de algumas regras de tipografia.

\section{Exemplo de código}
%%%%%%%%%%%%%%%%%%%%%%%%%%%%%%%%%%%%%%%%

\begin{lstlisting}[numbers=none,language=java,caption={[CommandDaemonCallsItf]
   {CommandDaemon} callback interfaces},label=lis:commandDCallsItfs,float=htb]
public interface CallBackCmdMeasurements { // comment
	public abstract void newMeasure(MeasurementBasic measure, int reqId);
	public abstract void newMeasuresAggSimp(MeasurementBasic[] measuresAggSimp, 'A string');
}
\end{lstlisting}

É possível como referir o código, por exemplo o bloco de código~\ref{lis:commandDCallsItfs}.

\section{Acrónimos}
%%%%%%%%%%%%%%%%%%%%%%%%%%%%%%%%%%%%%%%%

Deve-se acrescentar os acrónimos no ficheiro \texttt{acros.tex} e ordená-los alfabeticamente nesse ficheiro.
Vamos usar o acrónimo \ac{TCP} que deve estar expandido, assim como no capítulo~\ref{chap:tests}. Os acrónimos devem aparecer expandidos em cada capítulo (o que está já configurado para esta dissertação).

Podem usar $\backslash$\texttt{acs} para apenas mostrar o acrónimo, $\backslash$\texttt{acl} para mostrar a expansão. Ver mais na documentação do pacote \texttt{acronym}.

\section{Figuras}
%%%%%%%%%%%%%%%%%%%%%%%%%%%%%%%%%%%%%%%%

Podem ver a figura~\ref{fig:logoFCUP} muito bem. Notem que na lista de figuras não aparece tudo o que está na legenda, mas apenas o que está entre~[].

\begin{figure}[htb]
   \centering % center the figure
   \includegraphics[scale=.4]{pics/fc_logo}
   \caption[FCUP logo velho]{O logo da FCUP antigo}\label{fig:logoFCUP}
\end{figure}

As figuras devem usar o posicionamento [htb] para serem preferencialmente colocadas no local onde foram declaradas.

\subsection{Formato}
%---------------------------------------
    Devem colocar as figuras sempre que possível em pdf. Usem o programa (ou aplicação web) para exportar para pdf. Desse modo a qualidade das mesmas é sempre maior.

    No caso em que o pdf gerado ocupe uma página A4 inteira (por exemplo nas ferramentas do MS Office) podem usar o \texttt{\href{https://www.ctan.org/pkg/pdfcrop?lang=en}{pdfcrop}} no pdf gerado. Este comando está disponível na maior das distribuições de \LaTeX.


\section{Referências a bibliografia}
%%%%%%%%%%%%%%%%%%%%%%%%%%%%%%%%%%%%%%%%
Algumas referências para se ver a ordenação~\cite{yaacoub2012} (estão ordenadas pelo apelido do autor). Aqui outra ainda~\cite{etsitr102732}.

Temos aqui o Clausen~\cite{Clausen2003}. E para ver várias~\cite{yaacoub2012, etsitr102732, strunk2007elements}


\section{Sugestões na escrita e uso do tex}
\begin{itemize}
  \item usar o \textasciitilde{} para ligar as referências, evita a possível mudança de página: ex.: como o Brandão refere\textasciitilde\textbackslash{}cite\{bran99\}

\end{itemize}
\section{Sugestões na escrita e uso do tex}
\begin{itemize}
  \item em inglês não usar a \emph{short form}. Em textos formais deve-se manter a \emph{long form}. Incorreto: ``don't use''. \textbf{Correto:} ``do not use''
  \item em inglês (e também português) não usar a forma reflexiva ou indireta, preferir sempre a forma direta (que é mais assertiva e torna as frases menos complexas).
     \begin{itemize}
        \item incorreto:``Python was used to program''; Correto: ``we used Python''
     \end{itemize}
\end{itemize}

\section{Tabelas}
%%%%%%%%%%%%%%%%%%%%%%%%%%%%%%%%%%%%%%%%
Exemplo de uma tabela mais complexa na tabela~\ref{tab:bsnvswsn}.

\begin{table}[htbp]
   \caption[BSN vs WSN]{BSN versus WSN
   (with input from  Latré  and Guang )}
\label{tab:bsnvswsn}
\centering
{
\footnotesize
\begin{tabularx}{0.98\textwidth}{|>{\columncolor{gray-cell}}c|X|X|}
   \hline
   \rowcolor{gray-cell} 
   &    \centering  \textbf{BSN} &  \centering \textbf{WSN}  \tabularnewline 
   \hline
   %%%% line
   \begin{sideways} \hspace{-11em} \textbf{Distribution} \end{sideways}
   &  
   \begin{asparaenum}[\bfseries i)]
      \item Existence of a \ac{BS};
      \item \ac{BS} collects, maintains and processes the data;
      \item Nodes will do minimal processing, sending all data to the \ac{BS};
      \item Centralized system where \ac{BS} controls all nodes;
      \item Node replacement is difficult in in-body sensor nodes;
      \item Smaller number of nodes;
      \item Nodes need to take biocompatibility, wearability into account.
   \end{asparaenum}
   & 
   \begin{asparaenum}[\bfseries i)]

      \item A \ac{BS} may or not exist or there may be several \acp{BS} (e.g. mobile nodes 
         collect info, clustering);
      \item As in \ac{BSN}, but also on-demand querying;
      \item Nodes will do processing, aggregation to alleviate communication or 
         correlate results;
      \item Distributed system, nodes decide cooperatively;
         %(form clusters, aggregate data, \etc);
      \item Node replacement is difficult due to location, scale, etc.;
      \item (usually) Wide areas covered by large number of nodes.
      \item Nodes may need to be environment friendly, indiscernible from surroundings.
     %    \vspace{-0.7em}
   \end{asparaenum}
   \tabularnewline \hline
   %%%% line
   \textbf{Comm.} 
   & 
   \begin{asparaenum}[\bfseries i)]
      \item One hop to \ac{BS};
      \item Close range but attenuated by body;
      \item Data rates heterogeneous.
    %\vspace{-0.7em}
   \end{asparaenum}
   & 
   \begin{asparaenum}[\bfseries i)]
      \item Multi hop through network of sensor nodes;
      \item Long(er) range;
      \item Data rates homogeneous.
    %\vspace{-0.7em}
   \end{asparaenum}
   \tabularnewline \hline
   %%%% line
    \textbf{Data} 
   & 
   This is some text on this cell. The multirow package does not know the height of the cell and can not center the cell to the right. This is because of the X from tabularx.
   & 
   \multirow{ 2}{*}{This is on two rows of the table}

 \tabularnewline \cline{1-2}\noalign{\vskip.3pt}% using the noaling vskip to show the line
   %%%% line
    \textbf{Energy} 
   & 
   Some more text just to show something
   %      \vspace{-0.7em}
   & 
    \tabularnewline \hline
\end{tabularx}
}
\end{table}

%%%%%%%%%%%%%%%%%%%%%%%%%%%%%%%%%%%%%%%%
%%%%%%%%%%%%%%%%%%%%%%%%%%%%%%%%%%%%%%%%
\section{Referências para outras fontes de informação}
%%%%%%%%%%%%%%%%%%%%%%%%%%%%%%%%%%%%%%%%



%\chapter{Desenvolvimento}\label{chap:dese}
\chapter{Experiências e Testes}\label{chap:tests}

O referente ao ``Materials and Methods'' do capítulo~\ref{chap:back}.

Outro acrónimo pode ser \ac{TCP} (que deve estar expandido aqui, apesar de ter sido usado já no capítulo \ref{chap:devel}).

\lipsum
 

%\chapter{Resultados e análise}\label{chap:results}
\chapter{Resultados e análise}\label{chap:results}
 

%\chapter{Conclusões}\label{chap:conc}
\chapter{Conclusões}\label{chap:conc}

Teste de citações e sua ordem~\cite{nielsen1994usability,nielsen1994usabilityBook, nielsen1993mathematical}

\lipsum[4-6]

\section{Trabalho Futuro}\label{sec:trab}

\lipsum[1-3]



%% appendix
\appendix
%% \include{app1}

%% references
%\renewcommand{\bibname}{Referências} % o babel portuguese coloca Bibliografia
% os meses do ficheiro bib poderão aparecer em inglês, caso se pretenda deve-se colocar o texto em português explicitamente no ficheiro bid
\cleardoublepage%
\phantomsection%
\addcontentsline{toc}{chapter}{\bibname}
\bibliographystyle{plainnaturlAuthor} % use plainnaturlAppear to order references by appearance 
% usually it is by author on thesis, to ease Author lookup
%\nocite{*}  % Include all entries in references.bib, not just the ones cited.
\bibliography{refs} %changed the env to make it a numbered chapter

%% bye
\end{document}
